\chapter{五族残骸}

避难点的空气系统撑了不到二十分钟。

它没有彻底坏掉，只是开始间歇性抽风。每一次抽风，墙缝里都会喷出带灰的冷气，把地上的尘和死者身上的碎屑吹起来。活着的人咳嗽，死去的人沉默。那种沉默比咳嗽更重。

洛安和瑟恩不得不清理出一条通向后门的路。

所谓后门，是避难点最里面的一块维修板。板后可能连着旧管道，也可能只连着塌方。没有人知道。可前门外的裂缝正在逼近，空气也撑不了多久，不知道的路总比确定的死路好。

清路意味着搬开尸体。

这件事比想象中更难。

日辉族士兵很重。他们骨架修长，死后身体仍保持某种仪式感，像连倒下都要符合规训。洛安拖动其中一具时，尸体额头上的金属环刮过地面，发出清脆声响。金属环上刻着细密线条，中心是太阳形符号。

一段记忆从洛安脑中浮起：日辉族军官在战前宣读律令，声音通过扩音塔传遍阵地。他们说法厄同必须被统一，地球必须被交给拥有秩序的族群，否则未来生命只会在混乱中互相吞噬。

洛安当时站在队列末尾，只记得那军官的头发在灯下亮得刺眼。

现在，那种亮埋在灰里。

鳞裔族女人更难搬。

问题不在重量。她死死抱着那只封存罐。洛安试图把罐子拿出来，尸体僵硬的手指却扣得极紧。瑟恩看了一眼，说：“别动她的东西。”

“我们需要挪开她。”

“那就连东西一起挪。”

洛安照做。

封存罐裂缝里有绿色微光，像某种还没完全死去的东西。鳞裔族的鳞片没有蜥蜴那样的粗糙感，冷硬、细密，边缘带着金属质感。她脖颈处挂着一枚黑色小牌，上面没有名字，只有一串周期编号。

洛安认得这种编号。

长期库。

鳞裔族不相信短暂胜利。他们保存基因、保存胚胎，也保存失败后的下一次机会。战争在他们眼里跨过一代人的寿命，可以被拖进地下、拖进睡眠、拖进漫长周期。

角落里的节肢族个体最轻，也最难让人下手。

它们像昆虫，却又明显不是昆虫。外骨骼呈暗褐色，胸部嵌着战术节点，复眼关闭后像两块碎裂黑镜。洛安搬起其中一具时，另一具的肢节突然动了一下。

避难点里所有人都僵住。

那具节肢族没有复活。只是某个残留指令让它的前肢抬起，轻轻碰了碰同伴的胸部节点。动作很小，像在确认连接是否还在。

然后它彻底不动了。

查理，也就是我，在那一瞬间忽然不那么害怕它们了。

如果一个集群意识的个体在死后仍试图确认同伴，那它们真的没有感情吗？还是我们只是听不懂它们表达感情的方式？

渊民族幸存者主动站起来，帮忙挪开自己族人的尸体。

那具渊民族战士死在门边，鳃片干裂，皮肤从蓝灰变成近乎白色。幸存者跪在他身旁，用随身水囊里最后一点水抹过死者额头。水一接触皮肤，便发出极低的声音。

不是滴水声。

像歌。

避难点里没人说话。

那声音很短，却让我想起巨鲸坠落前的低频。不同的是，巨鲸的声音像星球骨骼被拉动，而这声音像海在很远的地方记住了一个名字。

“渊民会把死者交给水。”洛安的记忆告诉我，“如果没有水，就交给声音。”

最后是震翼族幼兵。

他还活着，但失去了半边翅膀和一条腿。瑟恩靠过去时，幼兵睁开眼，银蓝色虹膜已经开始浑浊。他用极弱的翅振说了些什么。

瑟恩的脸色变了。

洛安听不懂全部，只捕捉到几个频率：

上层通道。

仍有人。

不要让日辉族拿到种子。

种子。

这个词像一枚小石子落进我意识里。

洛安也听见了，但他不知道意思。瑟恩显然知道一点，却没有立刻解释。他只把幼兵的手放回胸前，低声说：“收到。”

幼兵的翅膀停止震动。

维修板终于被清出来。

洛安回头看了一眼避难点。短短二十分钟里，这里像一座被压缩的太阳系坟场。日辉族、鳞裔族、节肢族、渊民族、震翼族，所有刚才还在战场上互相杀戮的族群，现在被迫挤在同一个快要失效的盒子里。

死亡把他们混在一起。

而活着的人不得不从他们中间爬出去。

瑟恩撬开维修板，一股冷风从后面涌出。

不是新鲜空气。

是通向更深处的灰。

洛安握紧短刃，钻了进去。

我跟着他，穿过五族残骸之间最后一点狭窄的空隙。
